% chapters/14/sub-cpmk.tex
\begin{subcpmkbox}
Bab ini membuka \CPMK{4}, yaitu capaian terakhir pada mata kuliah ini. Yang
dibahas bukan lagi mekanisme permainan, melainkan penyajiannya kepada pemain.
Bab ini menuntaskan \SUBCPMK{4.1} secara utuh.

\vspace{0.5em}
\textbf{\SUBCPMK{4.1} Mengimpor dan mengonfigurasi aset visual
(\istilah{sprite}/model, animasi, UI) ke dalam \istilah{scene}} --- ranah C3,
\CPMK{4}, \CPL{04}.

\vspace{0.5em}
{\setlength{\tabcolsep}{4pt}\footnotesize%
\begin{tabularx}{\textwidth}{@{}C{1.9cm} L{4.0cm} L{2.5cm} Y@{}}
\toprule
\headrow \textbf{Indikator} & \textbf{Rumusan} & \textbf{Dibahas pada} & \textbf{Bukti} \\
\midrule
\ind{a} & Mengimpor aset dan mengatur \istilah{importer setting} dengan tepat &
Bagian~\ref{sec:pipeline-impor} &
Catatan pengaturan impor yang dipakai pada aset sendiri, beserta alasan setiap
pilihannya. \\
\addlinespace
\ind{b} & Menyiapkan \istilah{Animator} beserta transisi animasi (misalnya
\istilah{idle} dan \istilah{run}) &
Bagian~\ref{sec:animator} &
Animasi diam dan berlari yang berjalan bergantian menurut keadaan pemain,
beserta tangkapan layar jendela \istilah{Animator}. \\
\addlinespace
\ind{c} & Menampilkan elemen UI (skor, nyawa) yang berfungsi pada
\istilah{scene} &
Bagian~\ref{sec:canvas-ui} &
Antarmuka yang menampilkan skor dan nyawa, beserta fungsinya yang terhubung
pada peristiwa permainan. \\
\bottomrule
\end{tabularx}}
\end{subcpmkbox}
